\section*{Physique (13 points)}
\section*{Exercice 1 ( 2,5 points) : Âge d'une nappe phréatique}
Le chlore existe dans la nature sous forme de trois isotopes~: le chlore 35
$(\ce{_{17}^{35}Cl})$, le chlore 36 $(\ce{_{17}^{36}Cl})$ et le chlore 37
$(\ce{_{17}^{37}Cl})$.

Dans les eaux de surface (mers, lacs, \ldots), le chlore 36 est constamment
renouvelé et sa teneur peut être supposée constante.

Dans le cas des eaux de nappes phréatiques, le renouvellement n'existe plus et
la proportion en chlore 36 diminue au cours du temps.

\begin{donnees}
    \begin{table}[H]
        \centering
        \setlength{\arrayrulewidth}{0.8pt}
        \renewcommand{\arraystretch}{1.5}
        \begin{tabularx}{\textwidth}{|p{4cm}|C|C|C|}
            \hline
            Noyau ou particule & Électron & Chlore $\ce{_{17}^{36}Cl}$ &
            Argon $\ce{_{18}^{36}Ar}$ \\
            \hline
            Masse en (\unit{\u}) & $0,000549$ & $35,968312$ & $35,967545$ \\
            \hline
            \multicolumn{3}{|l|}{Constante radioactive du chlore 36~:
            $\lambda=\qty{2.30e-6}{\per\mathrm{ans}}$} &
            \multicolumn{1}{c|}{$\qty{1}{\u}=\qty{931.5}{\U}$} \\
            \hline
        \end{tabularx}
        \begin{tabularx}{\textwidth}{|p{10cm}|C|C|C|}
            \hline
            Noyau & $\ce{_{17}^{35}Cl}$ & $\ce{_{17}^{36}Cl}$ & 
            $\ce{_{17}^{37}Cl}$ \\
            \hline
            \raisebox{0mm}[7mm][4mm]{
                Énergie de liaison par nucléon~:
                $\frac{E_{\ell}}{A}$ ($\unit{\MeV\per\mathrm{nucléon}}$)} &
            $8,5178$ & $8,5196$ & $8,5680$ \\
            \hline
        \end{tabularx}
    \end{table}
\end{donnees}

\begin{enumerate}
    \item Recopier sur votre copie le numéro de la question et écrire la lettre
          correspondante à la proposition vraie.

          La composition du noyau de chlore $\ce{_{17}^{35}Cl}$ est~:

          \begin{minipage}{\linewidth}
              \begin{table}[H]
                  \centering
                  \setlength{\arrayrulewidth}{0.8pt}
                  \renewcommand{\arraystretch}{1.5}
                  \begin{tabularx}{\textwidth}{|p{1cm}|C|}
                      \hline
                      A & 17 protons et 35 neutrons \\
                      \hline
                      B & 18 protons et 17 neutrons \\
                      \hline
                      C & 17 protons et 18 neutrons \\
                      \hline
                      D & 18 protons et 35 neutrons \\
                      \hline
                  \end{tabularx}
              \end{table}
          \end{minipage}
    
    \item Déterminer, en justifiant votre réponse, le noyau le plus stable parmi
          $\ce{_{17}^{35}Cl}$, $\ce{_{17}^{36}Cl}$ et $\ce{_{17}^{37}Cl}$.

    \item Le chlore 36 radioactif, donne en se désintégrant un noyau d'argon
          $\ce{_{18}^{36}Ar}$.
          \begin{enumerate}
              \item Écrire l'équation de désintégration d'un noyau de chlore 36
                    et identifier le type de cette désintégration.
              \item Calculer, en unité (\unit{\MeV}), l'énergie libérée
                    $E_{\text{libérée}}=|\Delta E|$ au cours de la
                    désintégration d'un noyau de chlore 36.
          \end{enumerate}

    \item Un échantillon de volume $V$, des eaux de surface, contient $N_{0}$
          noyaux de chlore 36. Un échantillon de même volume $V$, d'eau issue
          d'une nappe phréatique ne contient que $38 \%$ du nombre de noyaux de
          chlore 36 trouvé dans les eaux de surface.
    
          Déterminer, en unité (ans), l'âge de la nappe phréatique.
\end{enumerate}



\section*{Exercice 2 ( 5,5 points): Dipôle RC - Circuit RLC série}
Le condensateur, la bobine et le conducteur ohmique sont des composants électroniques dont le comportement diffère selon les circuits électriques ou ils se trouvent. Le condensateur et la bobine constituent des réservoirs d'énergie alors que le conducteur ohmique joue un rôle différent en agissant sur le bilan énergétique dans ces circuits.\\
Cet exercice vise :

\section*{- l'étude de la charge d'un condensateur;}
\begin{itemize}
  \item l'étude des oscillations électriques libres dans un circuit RLC série.\\
Le montage de la figure (1) comporte un générateur de tension de force électromotrice $E$, un conducteur ohmique de résistance $R$ réglable, un condensateur de capacité $C$, une bobine ( $L ; r$ ) et deux interrupteurs $K_{1}$ et $K_{2}$.
\end{itemize}

\begin{enumerate}
  \item On règle la résistance sur la valeur $R=100 \Omega$ et on ferme\\
\includegraphics[width=\textwidth]{2025_03_31_52b4902f5779608975f8g-4}
\end{enumerate}

Figure (1) $K_{1}$ à $t_{0}=0$, en maintenant $K_{2}$ ouvert.\\
1.1. Établir l'équation différentielle vérifiée par la tension $u_{C}(t)$ aux bornes du condensateur.\\
1.2. Un système d'acquisition permet d'obtenir les courbes de la figure (2) qui représentent $u_{C}(t)$ et $u_{R}(t)$ la tension aux bornes du conducteur ohmique.\\
1.2.1. Identifier la courbe correspondante à $u_{C}(t)$.\\
1.2.2. Déterminer graphiquement la valeur de :\\
a. la constante de temps $\tau$.\\
b. la force électromotrice $E$.\\
1.2.3. Vérifier que $C=50 \mu F$.\\
1.2.4. Déterminer la valeur maximale $I_{0}$ de l'intensité du courant électrique qui traverse le circuit.\\
1.2.5. La solution de l'équation différentielle demandée dans la question (1.1.) s'écrit :\\
$u_{C}(t)=E .\left(1-e^{-\frac{t}{\tau}}\right)$\\
\includegraphics[width=\textwidth]{2025_03_31_52b4902f5779608975f8g-4(1)}

Figure (2)

Recopier sur votre copie le numéro de la question et écrire la lettre correspondante à la proposition vraie.\\
L'expression de l'instantsité $i(t)$ du courant électrique qui traverse le circuit est :

\begin{center}
\begin{tabular}{|l|l|l|l|l|l|l|l|}
\hline
$\mathbf{A}$ & $i(t)=0,1 . e^{-200 . t}$ & $\mathbf{B}$ & $i(t)=0,1 . e^{-\frac{1}{200}}$ & $\mathbf{C}$ & $i(t)=0,1 .\left(1-e^{-200 \prime}\right)$ & $\mathbf{D}$ & $i(t)=0,1 . e^{-10, t}$ \\
\hline
\end{tabular}
\end{center}

0,25\\
2. Le condensateur étant chargé, on ouvre $K_{1}$ et on ferme $K_{2}$ à l'instant ( $t_{0}=0$ ).

À l'aide du même système d'acquisition on obtient la courbe de la figure (3) qui représente $u_{C}(t)$.\\
$\mathbf{0 , 2 5}$ 2.1. Nommer le régime d'oscillations que montre la courbe de la figure (3).\\
2.2. Déterminer la valeur de l'inductance $L$. On suppose que la pseudo période $T$ est égale à la période propre des oscillations libres du circuit ( $L C$ ). On prend ( $\pi^{2}=10$ ).\\
2.3. On note respectivement $\mathcal{E}_{e 0}$ et $\mathcal{E}_{\ell 1}$ les énergies électriques emmagasinées dans le condensateur aux instants $t_{0}=0$ et $t_{1}=T$.\\
0,5 2.3.1. Déterminer les valeurs de $\mathcal{E}_{e 0}$ et $\mathcal{E}_{e 1}$.\\
0,5\\
2.3.2. Calculer $\Delta \mathcal{E}$ la variation de l'énergie totale du circuit entre $t_{0}=0$ et $t_{1}=T$. Expliquer ce résultat.\\
\includegraphics[width=\textwidth]{2025_03_31_52b4902f5779608975f8g-5}

Figure (3)

\input{physique_ex3}



\section{Étude de mouvement d'un skieur - Étude d'un système oscillant}
Les deux parties 1 et 2 sont indépendantes

Les mouvements rectiligne, plan et oscillatoire sont des mouvements de types différents. Ces mouvements dépendent de la nature des milieux où ils se produisent, des types d'actions mécaniques appliquées et des conditions initiales.\\
Cet exercice vise:

\begin{itemize}
  \item l'étude de mouvement d'un skieur soumis à des forces constantes;
  \item l'étude de mouvement d'un solide soumis à une force variable.
\end{itemize}

\subsection{Étude de mouvement d'un skieur}
Un skieur aborde une piste horizontale AB. On modélise le skieur avec ses accessoires par un solide ( $S$ ), de masse $m$ et de centre d'inertie $G$.

\begin{enumerate}
  \item Le mouvement du solide ( $S$ ) sur la piste AB se fait avec frottement équivalent à une force unique $\bar{f}$ constante et de sens opposé au vecteur vitesse du skieur. Pour étudier le mouvement de $(S)$ sur le trajet AB, on choisit un repère $(O, \vec{i})$ lié à la Terre supposé galiléen,\\
\includegraphics[width=\textwidth]{2025_03_31_52b4902f5779608975f8g-5(1)}
\end{enumerate}

Figure (1) et l'instant de passage de $G$ en A comme origine des dates $\left(t_{0}=0\right)$. On repère la position de $G$ à un instant $t$ par son abscisse x dans ce repère. À $t_{0}=0: x_{G}=x_{0}=0$ (figure 1).\\
Données: $f=70 \mathrm{~N} ; m=70 \mathrm{~kg} ; g=10 \mathrm{~m} . \mathrm{s}^{-2}$\\
0,75 1.1. En appliquant la deuxième loi de Newton, établir l'équation différentielle vérifiée par l'abscisse $x_{G}$\\
0,5 1.2. Déterminer la nature du mouvement de $G$. Calculer l'accélération $a_{G}$ du mouvement de $G$.\\
0,5 1.3. Le skieur passe en A avec la vitesse $V_{A}=25 \mathrm{~m} \cdot \mathrm{~s}^{-1}$ et parcourt le trajet AB pendant une durée égale à $4,4 \mathrm{~s}$. Montrer que le skieur ne peut éviter la chute après la position B .\\
2. Le skieur passe en B avec une vitesse horizontale $\vec{V}_{B}$. Il tombe en chute libre sur le sol situé à la hauteur $h=B C=3,2 m$ de la piste $A B$ et touche le sol en un point $P$ d'abscisse $x_{P}=16,48 \mathrm{~m}$ dans le repère orthonormé ( $B, \vec{i}, \vec{j}$ ) lié à la Terre supposé galiléen. On choisit comme nouvelle origine des dates, l'instant de passage de $G$ en $B$.\\
Les équations horaires du mouvement de $G$ s'écrivent: $x_{G}(t)=V_{B} . t$ et $y_{G}(t)=\frac{1}{2} g . t^{2}$\\
0,5 2.1. Déterminer l'instant $t_{P}$ où le skieur touche le sol au point $P$.\\
0,5 2.2. Pour améliorer sa performance, le skieur a réalisé un deuxième essai sur la même piste AB . Il est passé en B avec une vitesse $V^{\prime}{ }_{B}$ pour atteindre une portée $x_{p}^{\prime}=18 \mathrm{~m}$.\\
Déterminer la valeur de la vitesse $V_{B}^{\prime}$.

\subsection{Étude d'un système oscillant}
\begin{minipage}[t][][t]{0.72\linewidth}
Un solide $(S)$ de masse $m$ est fixé à un ressort horizontal à spires non
jointives, de masse négligeable et de raideur $K$. À l'équilibre, le centre
d'inertie $G$ de $(S)$ coïncide avec l'origine du repère $(O,\vec{i})$ lié à la
Terre supposé galiléen.

On écarte ($S$) de sa position d'équilibre d'une distance $X_m$ et on
l'abandonne sans vitesse initiale. L'équation horaire du mouvement de $G$
est~: \\
$x(t)=X_m \cdot \cos\left(\frac{2\pi}{T_0} \cdot t+\varphi\right)$.
\end{minipage}
\hfill
\begin{minipage}[t][][t]{0.26\linewidth}
\begin{figure}[H]
    \flushright
    \begin{tikzpicture}
    \node[inner sep=0pt,minimum width=4.5cm,minimum height=3mm,
        top color=black!70] (sol) {};
    % Axe xx'
    \coordinate (O) at ($(sol.south west)!.8!(sol.south east)$);
    \draw[-Latex] (sol.south west)
        node[below right,inner xsep=0pt] {$x^{\prime}$} -- (O)
        node[below=0.4mm] {$O$} -- (sol.south east) -- ++(0.7,0)
        node[below left,inner xsep=0pt,inner ysep=2mm] {$x$};
    % Fixed
    \node[inner sep=0pt,draw,minimum width=3mm,minimum height=7mm,
        anchor=south,pattern=north west lines,
        thick] (fix) at ([xshift=3mm,yshift=0.4pt]sol.north west) {};
    % Solide (S)
    \node[inner sep=0pt,draw,minimum width=7mm,minimum height=7mm,
        anchor=south,rounded corners=2pt,thick,label=$\mathrm{(S)}$,
        fill=lightgray] (S) at
        ([yshift=0.4pt]O |- sol.north) {};
        %([xshift=-6mm,yshift=0.4pt]sol.north east) {};
    \node[inner sep=1pt,circle,fill,
        label={[right=-0.8mm,font=\scriptsize]:$G$}] (G) at (S.center) {};
    \draw[{Rectangle[length=1pt,width=4pt]}-{Stealth[length=5pt]},
        thick] (O) -- ++(0.8,0)
        node[below left,inner xsep=0pt] {$\ex$};
    % Ressort
    \draw[decoration={
          coil,
          aspect=0.3,
          segment length=2.0mm,
          amplitude=2.4mm,
          pre length=1mm,
          post length=1mm},
          decorate,thick,black] (fix.east) -- (S.west);
\end{tikzpicture}

    \caption{}
    \label{fig:solide_ressort}
\end{figure}
\end{minipage}

\begin{donnees}
    \begin{itemize}
        \item Tous les frottements sont négligeables~;
        \item $m=\qty{255}{\gram}$.
    \end{itemize}
\end{donnees}
\begin{enumerate}
    \item L'équation de la vitesse de $G$ s'écrit~:
          $v(t)=-0,25 \cdot \sin(2\pi\cdot t)$ \; (\unit{\v}).
          \begin{enumerate}
              \item En exploitant l'équation de la vitesse, déterminer la
                    période propre $T_0$ des oscillations, la valeur de
                    l'amplitude $X_m$ et la phase $\varphi$ à $t_0=0$.
              \item Vérifier que la raideur du ressort est
                    $K\approx\qty{10}{\newton\per\meter}$.
          \end{enumerate}
    \item Déterminer l'expression de la force de rappel $\vec{F}$ exercée par le
          ressort sur le solide ($S$) à l'instant $t=\qty{0,5}{\second}$.
\end{enumerate}
